\documentclass[12pt,a4paper]{article}

% ==========================================
% 1. POLICES ET ENCODAGE DE TEXTE
% ==========================================
\usepackage[utf8]{inputenc}
\usepackage[T1]{fontenc}        % Encodage pour les caractères accentués
\usepackage{mathptmx}          % Police Times harmonisée pour le texte et les formules mathématiques
\usepackage{pifont}            % Symboles graphiques (ex: \ding{45})

% ==========================================
% 2. DISPOSITION DE LA PAGE ET EN-TÊTES
% ==========================================
\usepackage[left=1cm, right=1cm, top=1cm, bottom=1.7cm]{geometry} % Marges du document
\usepackage{fancyhdr}          % Gestion des en-têtes et pieds de page
\usepackage{lastpage}          % Référence dynamique à la dernière page pour "Page X/Y"
\usepackage{multicol}          % Support pour la mise en page multi-colonnes

% --- Configuration de Fancyhdr ---
\pagestyle{fancy}
\fancyhf{} % Efface les configurations par défaut
\renewcommand{\headrulewidth}{0pt} % Suppression de la ligne d'en-tête
\renewcommand{\footrulewidth}{1pt} % Ligne de séparation de pied de page
\renewcommand{\footruleskip}{10pt} % Espacement vertical au-dessus du pied de page

\fancyfoot[L]{\color{blue}\ding{45}\ \textbf{Prof: M. BA}}
\fancyfoot[R]{\color{blue}\ding{45}\ \textbf{\the\year}} % Mise à jour automatique de l'année
\cfoot{\textbf{\thepage\ / \pageref{LastPage}}}

% ==========================================
% 3. MATHÉMATIQUES
% ==========================================
\usepackage{amsmath}           % Environnements et commandes mathématiques fondamentales
\usepackage{amssymb}           % Symboles mathématiques avancés
\usepackage{mathrsfs}          % Lettres calligraphiques via \mathscr{}

% Passage automatique de toutes les formules en mode affiché (displaystyle)
\everymath{\displaystyle}

% ==========================================
% 4. GESTION DES TABLEAUX ET OBJETS FLOTTANTS
% ==========================================
\usepackage{array}             % Extensions des capacités des tableaux
\usepackage{multirow}          % Fusion horizontale et verticale de cellules
\usepackage{varwidth}          % Boîtes à largeur adaptable
\usepackage{float}             % Placement strict des objets flottants via l'option [H]

% ==========================================
% 5. LISTES PERSONNALISÉES
% ==========================================
\usepackage{enumitem}          % Personnalisation fine des environnements d'énumération

% Définitions des puces personnalisées TikZ
\newcommand\circitem[1]{%
  \tikz[baseline=(char.base)]{%
    \node[circle, draw=gray, fill=red!55, minimum size=1.2em, inner sep=0pt] (char) {#1};%
  }%
}
\newcommand\boxitem[1]{%
  \tikz[baseline=(char.base)]{%
    \node[fill=cyan, minimum size=1.2em, inner sep=0pt] (char) {#1};%
  }%
}

\setlist[enumerate,1]{label=\protect\circitem{\arabic*}}
\setlist[enumerate,2]{label=\protect\boxitem{\alph*}}

% ==========================================
% 6. GRAPHISMES ET TABLEAUX DE VARIATIONS
% ==========================================
\usepackage{tikz}
\usetikzlibrary{patterns, trees, calc, babel} % Regroupement des extensions TikZ
\usepackage{tkz-tab}           % Conception de tableaux de variations et de signes

% ==========================================
% 7. ENCADRÉS STYLISÉS (TCOLORBOX)
% ==========================================
\usepackage[most]{tcolorbox}
\tcbuselibrary{skins, hooks}

% Environnement d'encadré d'exemple
\newtcolorbox{exa}[2][]{%
  enhanced,
  breakable,
  before skip=2mm,
  after skip=5mm,
  colback=yellow!20!white,
  colframe=black!20!blue,
  boxrule=0.5mm,
  attach boxed title to top left={xshift=0.6cm, yshift*=1mm-\tcboxedtitleheight},
  fonttitle=\bfseries,
  title={#2},
  #1,
  boxed title style={%
    frame code={%
      \path[fill=tcbcolback!30!black]
        ([yshift=-1mm,xshift=-1mm]frame.north west)
        arc[start angle=0,end angle=180,radius=1mm]
        ([yshift=-1mm,xshift=1mm]frame.north east)
        arc[start angle=180,end angle=0,radius=1mm];
      \path[left color=tcbcolback!60!black,right color=tcbcolback!60!black,middle color=tcbcolback!80!black]
        ([xshift=-2mm]frame.north west) -- ([xshift=2mm]frame.north east)
        [rounded corners=1mm]-- ([xshift=1mm,yshift=-1mm]frame.north east)
        -- (frame.south east) -- (frame.south west)
        -- ([xshift=-1mm,yshift=-1mm]frame.north west)
        [sharp corners]-- cycle;
    },
    interior engine=empty,
  },
  interior style={top color=yellow!5}
}

% ==========================================
% 8. FILIGRANE (ARRIÈRE-PLAN)
% ==========================================
\usepackage{eso-pic}
\AddToShipoutPicture{
  \AtTextCenter{%
    \makebox[0pt]{\rotatebox{80}{\textcolor[gray]{0.7}{\fontsize{5cm}{5cm}\selectfont PGB}}}
  }
}

% ==========================================
% 9. COMPTEURS ET MACROS PEDAGOGIQUES
% ==========================================
\newcounter{exercice}
\renewcommand{\theexercice}{\arabic{exercice}}
\newcommand{\exo}{\refstepcounter{exercice}\textbf{Exercice \theexercice}\ }

% ==========================================
% 10. HYPERLIENS ET INTERACTIVITÉ
% ==========================================
\usepackage[colorlinks=true, linkcolor=blue, urlcolor=blue, citecolor=blue]{hyperref}

\begin{document}
\renewcommand{\arraystretch}{1.5}
\renewcommand{\arrayrulewidth}{1.2pt}
\begin{tikzpicture}[overlay,remember picture]
    \node[draw=blue,line width=1.2pt,fill=purple,text=blue,inner sep=3mm,rounded corners,pattern=dots]at ([yshift=-2.5cm]current page.north) {\begingroup\setlength{\fboxsep}{0pt}\colorbox{white}{\begin{tabular}{|*1{>{\centering \arraybackslash}p{0.28\textwidth}} |*2{>{\centering \arraybackslash}p{0.2\textwidth}|} *1{>{\centering \arraybackslash}p{0.19\textwidth}|} }
                \hline
                \multicolumn{3}{|c|}{$\diamond$$\diamond$$\diamond$\ \textbf{Lycée de Dindéfélo}\ $\diamond$$\diamond$$\diamond$ } & \textbf{A.S. : 2025/2026}                                              \\ \hline
                \textbf{Matière: Mathématiques}                                                                     & \textbf{Niveau : TS2} & \textbf{Date: 29/12/2025} & \textbf{} \\ \hline
                \multicolumn{4}{|c|}{\parbox[c]{10cm}{\begin{center}
                                                                  \textbf{{\Large TD : Fonction Exponentielle}}
                                                              \end{center}}}                                                                                                        \\ \hline
            \end{tabular}}\endgroup};
\end{tikzpicture}
\vspace{3cm}

% ==========================================
% EXERCICE 1 : CALCUL DE LIMITES
% ==========================================
\begin{center}
\fbox{\textbf{\exo}} Déterminer les limites suivantes :
\end{center}

\begin{multicols}{2}
\setlength{\columnseprule}{0.1mm}
\begin{enumerate}
\item $\lim_{x \to +\infty} e^{x+2}$
\item $\lim_{x \to -\infty} e^{3x-1}$
\item $\lim_{x \to +\infty} (e^x)^2$
\item $\lim_{x \to -\infty} \frac{1}{e^x}$
\item $\lim_{x \to +\infty} e^{x^2 + x + 1}$

\item $\lim_{x \to -\infty} e^{x^2 + x + 1}$
\item $\lim_{x \to +\infty} e^{\frac{x+1}{x^2 + x + 1}}$
\item $\lim_{x \to -\infty} e^{\frac{2x+3}{x-1}}$
\item $\lim_{x \to +\infty} \frac{e^{x+2}}{e^{x+1}}$
\item $\lim_{x \to -\infty} \frac{e^x + 2}{e^x + 1}$

\item $\lim_{x \to +\infty} \frac{e^x}{\sqrt{x}}$
\item $\lim_{x \to +\infty} \frac{e^x}{x^2}$
\item $\lim_{x \to +\infty} (e^x - \sqrt{x})$
\item $\lim_{x \to -\infty} x^2 e^x$
\item $\lim_{x \to -\infty} \sqrt{-x} \, e^x$

\item $\lim_{x \to +\infty} (e^x - 2x)$
\item $\lim_{x \to +\infty} \frac{e^{x^2 + 1}}{e^{x+1}}$
\item $\lim_{x \to 0} \frac{e^{2x} - 1}{x}$
\item $\lim_{x \to 0} \frac{e^{x^2} - 1}{x}$
\item $\lim_{x \to 0} \frac{e^{3x} - 1}{2x}$

\item $\lim_{x \to 0} \frac{e^{\sin x} - 1}{x}$
\item $\lim_{x \to -\infty} \frac{e^x - 1}{x}$
\item $\lim_{x \to +\infty} \frac{e^x - 1}{x}$
\item $\lim_{x \to 0} \frac{e^{x^2} - 1}{3x}$
\item $\lim_{x \to 0} \frac{e^{2x} - e^x}{x}$

\item $\lim_{x \to -\infty} \frac{2e^x + 1}{e^x + 3}$
\item $\lim_{x \to +\infty} \frac{2e^x + 1}{e^x + 3}$
\item $\lim_{x \to -\infty} \frac{e^x - 1}{2x}$
\item $\lim_{x \to +\infty} \frac{e^x - 1}{2x}$
\item $\lim_{x \to 2} e^{\frac{1}{2-x}}$

\item $\lim_{x \to +\infty} e^{x^2 - 2x}$
\item $\lim_{x \to -\infty} e^{x^2 + 4x}$
\item $\lim_{x \to 0} e^{\frac{2x - 1}{x - 2}}$
\item $\lim_{x \to 0^+} e^{\frac{1 + x}{x}}$

\item $\lim_{x \to +\infty} e^{\frac{1 - x}{-x - 1}}$
\item $\lim_{x \to -\infty} e^{\frac{x^3 + x}{-2x^2 - 1}}$
\item $\lim_{x \to 0} e^{|x^2 + x|}$
\item $\lim_{x \to 1} e^{|x^2 - 1|}$

\item $\lim_{x \to +\infty} e^{|x^2 - 1|}$
\item $\lim_{x \to +\infty} \frac{2x - 1}{x e^x - x}$
\item $\lim_{x \to -\infty} \frac{2x - 1}{x e^x + x}$
\item $\lim_{x \to +\infty} \frac{-1}{x^2 e^x}$

\item $\lim_{x \to -\infty} \frac{2e^x - 3}{3e^x + 1}$
\item $\lim_{x \to +\infty} \frac{2e^x - 3}{3e^x + 1}$
\item $\lim_{x \to -\infty} \frac{e^x - 3x}{2e^x + x}$
\item $\lim_{x \to +\infty} \frac{e^x - 3x}{2e^x + x}$

\item $\lim_{x \to +\infty} (x - e^x)$
\item $\lim_{x \to +\infty} \frac{1 - e^x}{x}$
\item $\lim_{x \to -\infty} \frac{1 - e^x}{x}$
\item $\lim_{x \to -\infty} (x e^x - x)$

\item $\lim_{x \to -\infty} x^n e^x \quad (n \in \mathbb{N}^*)$
\item $\lim_{x \to -\infty} x(e^x - 1)$
\item $\lim_{x \to +\infty} \frac{e^{2x}}{x}$
\item $\lim_{x \to +\infty} \frac{(x - 1)e^x}{x}$

\item $\lim_{x \to -\infty} \frac{(x - 1)e^x}{x}$
\item $\lim_{x \to +\infty} x \left(e^{\frac{1}{x}} - 1\right)$
\item $\lim_{x \to +\infty} x \left(e^{\frac{2}{x}} - 1\right)$
\item $\lim_{x \to -\infty} \frac{e^{-x}}{x}$

\item $\lim_{x \to -\infty} \left(e^x + \frac{1}{e^x}\right)$
\item $\lim_{x \to +\infty} \left( \frac{x^2}{x-1} - e^x \right)$
\item $\lim_{x \to -\infty} x^2 e^x$
\item $\lim_{x \to 0} \frac{e^x - 1}{x^2 + x}$

\item $\lim_{x \to 0} \frac{e^x - e^2}{x-2}$
\item $\lim_{x \to 0} \frac{e^{a+x} - e^a}{x} \quad (a \in \mathbb{R})$
\item $\lim_{x \to -\infty} \frac{e^x - 2}{e^x + 1}$
\item $\lim_{x \to 0} \tan x \cdot e^{\frac{1}{x}}$

\item $\lim_{x \to \frac{\pi}{4}} \frac{e^{\tan x} - e}{\sin x - \cos x}$
\item $\lim_{x \to 0} \frac{e^{\cos x - 1} - 1}{x^2}$
\item $\lim_{x \to +\infty} \frac{x e^x}{x+1}$
\item $\lim_{x \to +\infty} \frac{e^{1+x}}{1+x^2}$

\item $\lim_{x \to -\infty} \frac{e^x}{1 - e^x}$
\item $\lim_{x \to +\infty} \frac{e^x}{1 - e^x}$
\item $\lim_{x \to 0} \frac{e^{1-\cos x} - 1}{x^2}$
\item $\lim_{x \to 1} \left( -x + e^{\frac{x}{x-1}} \right)$

\item $\lim_{x \to -\infty} \left( -x + e^{\frac{x}{x-1}} \right)$
\item $\lim_{x \to +\infty} \left( -x + e^{\frac{x}{x-1}} \right)$
\item $\lim_{x \to +\infty} x \left( e^{\frac{x+1}{x-1}} - e \right)$
\item $\lim_{x \to +\infty} \frac{e^{\sqrt{x^2 - 1}}}{x^2 - 1}$

\item $\lim_{x \to 1} \frac{e^{(x-1)^2} - 1}{(x-1)^2}$
\item $\lim_{x \to 0} \frac{e^{\sqrt{x+1}-1} - 1}{x}$
\item $\lim_{x \to +\infty} \left( e^x + \frac{1}{e^x} \right)$
\item $\lim_{x \to -\infty} \frac{e^x}{1+x^2}$

\item $\lim_{x \to +\infty} \frac{e^x}{1+x^2}$
\item $\lim_{x \to +\infty} \frac{x^3}{e^x}$
\item $\lim_{x \to +\infty} \frac{e^{1+e^x}}{e^x}$
\item $\lim_{x \to +\infty} (x e^x - e^x)$

\item $\lim_{x \to +\infty} \sqrt{x} \, e^{-x}$
\item $\lim_{x \to 0} \frac{x}{e^{\sin x} - 1}$
\item $\lim_{x \to 0} \frac{e^{2x} - 1}{\tan x}$
\item $\lim_{x \to +\infty} \frac{e^x}{x-1}$
\item $\lim_{x \to 1} \frac{e^{x-1} - 1}{x-1}$
\end{enumerate}
\end{multicols}

% ==========================================
% EXERCICE 2 : ÉQUATIONS ET INÉQUATIONS
% ==========================================
\begin{center}
\fbox{\textbf{\exo}} Équations et inéquations dans $\mathbb{R}$ :
\end{center}

\begin{multicols}{2}
\setlength{\columnseprule}{0.1mm}
\textbf{1. Résoudre dans $\mathbb{R}$ les équations :}
\begin{enumerate}
\item $e^{2x} - 3e^x + 2 = 0$
\item $e^x - 2e^{-x} - 1 = 0$
\item $e^{3x-1} = 1$
\item $e^{x^2 - 4} = e^{3x}$
\item $\frac{e^x + 3}{e^x - 1} = 2$
\item $e^{2x} + e^x - 6 = 0$
\end{enumerate}

\columnbreak

\textbf{2. Résoudre dans $\mathbb{R}$ les inéquations :}
\begin{enumerate}
\item $e^{2x} - 4 > 0$
\item $e^x + 1 - 2e^{-x} \le 0$
\item $e^{x^2 - 1} < e^{x+1}$
\item $\frac{e^x - 1}{e^x + 2} \ge 0$
\item $e^{3x} - 2e^{2x} - e^x + 2 > 0$
\item $\frac{2}{e^x - 1} < 1$
\end{enumerate}
\end{multicols}

% ==========================================
% EXERCICE 3 : ÉTUDE DE FONCTIONS EXPONENTIELLES
% ==========================================
\begin{center}
\fbox{\textbf{\exo}} Étude de fonctions comportant la fonction exponentielle :
\end{center}

\begin{enumerate}[label=\arabic*.]
\item Soit $f$ la fonction définie sur $\mathbb{R}$ par : $f(x) = (x-1)e^x + 1$.
   \begin{enumerate}
   \item Étudier les limites de $f$ en $-\infty$ et $+\infty$.
   \item Étudier les variations de $f$ et dresser son tableau de variation.
   \item Déduire le signe de $f(x)$ sur $\mathbb{R}$.
   \end{enumerate}

\item Soit $g$ la fonction définie par : $g(x) = \frac{e^x - 1}{e^x + 1}$.
   \begin{enumerate}
   \item Déterminer l'ensemble de définition $D_g$ de $g$.
   \item Montrer que $g$ est une fonction impaire.
   \item Calculer la limite de $g$ en $+\infty$ et interpréter géométriquement le résultat.
   \item Calculer $g'(x)$ et dresser le tableau de variation de $g$.
   \item Déterminer l'équation de la tangente $(T)$ à la courbe de $g$ au point d'abscisse $0$.
   \end{enumerate}
\end{enumerate}

\end{document}